\chapter{Circle and Trigonometry}
\label{ch:rational-circle-trigonometry}

The purpose of this chapter is to construct trigonometric functions as
algorithms from the geometry of the unit circle, and then prove their
identities from that construction.  The chapter has four layers: rational
circle geometry, Archimedes-style exhaustion estimates, trigonometric
function algorithms, and the identities proved from those algorithms.

\section{Rational Circle}

\begin{definition}[Rational circle parametrization]
\label{def:rational-circle-point}
\lean{ComputableAnalysis.RationalCircle.Stage.point}
For a rational slope coordinate \(u\in\Q\), define
\[
  P(u)=
  \left(
    \frac{1-u^2}{1+u^2},
    \frac{2u}{1+u^2}
  \right).
\]
This is the rational stereographic chart for the unit circle, with the
missing point handled separately when needed.
\end{definition}

\begin{theorem}[Rational circle equation]
\label{thm:rational-circle-equation}
\lean{ComputableAnalysis.RationalCircle.Stage.point_normSq}
\uses{def:rational-circle-point}
\leanok
For every \(u\in\Q\),
\[
  P(u)_x^2+P(u)_y^2=1.
\]
\end{theorem}

\begin{definition}[Dyadic rational circle stages]
\label{def:dyadic-rational-circle-stages}
\lean{ComputableAnalysis.RationalCircle.dyadicStage}
\uses{def:rational-circle-point}
At stage \(n\), let \(N_n=2^n\).  For \(0\le k\le N_n\), the finite circle
sample is
\[
  C_{n,k}=P\left(\frac{k}{N_n}\right).
\]
The list \(C_{n,0},\ldots,C_{n,N_n}\) is the rational first-quadrant
sampling used by both the exhaustion estimates and the geometric point
algorithm.
\end{definition}

\begin{theorem}[Dyadic refinement]
\label{thm:dyadic-rational-circle-refinement}
\lean{ComputableAnalysis.RationalCircle.Trigonometry.dyadic_point_refineIndex}
\uses{def:dyadic-rational-circle-stages}
\leanok
The next dyadic stage keeps every old sample:
\[
  C_{n+1,2k}=C_{n,k}.
\]
The odd indices \(C_{n+1,2k+1}\) are the new rational points inserted between
old neighbors.
\end{theorem}

\begin{theorem}[Dyadic samples are circle samples]
\label{thm:dyadic-trig-samples}
\lean{ComputableAnalysis.RationalCircle.Trigonometry.dyadic_cos_sq_add_sin_sq}
\uses{def:dyadic-rational-circle-stages,thm:rational-circle-equation}
\leanok
For every stage \(n\) and index \(k\), the dyadic sample coordinates satisfy
\[
  C_{n,k,x}^2+C_{n,k,y}^2=1.
\]
The endpoints are
\[
  C_{n,0}=(1,0),
  \qquad
  C_{n,N_n}=(0,1).
\]
\end{theorem}

\section{Archimedes Method of Exhaustion}

\begin{definition}[Area and circumference enclosures]
\label{def:archimedes-area-circumference-enclosures}
\lean{ComputableAnalysis.piCircleArea}
\lean{ComputableAnalysis.piCircumference}
\lean{ComputableAnalysis.circumferenceSqrtPrecision}
\lean{ComputableAnalysis.RationalCircle.Stage.areaInterval}
\lean{ComputableAnalysis.RationalCircle.Stage.lengthPrecision}
\lean{ComputableAnalysis.RationalCircle.Stage.circumferenceInterval}
\uses{def:dyadic-rational-circle-stages}
\leanok
At each rational circle stage, the inner polygon and outer tangent polygon
give a rational area interval for the unit disk and a rational path-length
interval for the unit circle.  The corresponding raw reals are the area
algorithm \(\code{ComputableAnalysis.piCircleArea}\) and the circumference
algorithm \(\code{ComputableAnalysis.piCircumference}\).
For circumference, the polygon subdivision stage and the square-root
subroutine stage are deliberately separated.  If the polygon has \(N\)
subdivisions, each rational segment length is computed by asking the
square-root subroutine for the finer stage
\[
  (N+1)^3.
\]
Thus the main geometric gap is controlled by the polygon refinement, while
the accumulated square-root wiggle is kept smaller by a cheap internal
precision schedule.
The same finite geometry is later reused for sector-area comparisons in the
angle-to-circle-point algorithm.
\end{definition}

\begin{theorem}[Finite Archimedes bounds]
\label{thm:finite-archimedes-bounds}
\lean{ComputableAnalysis.PiProofs.FiniteArchimedesStage}
\lean{ComputableAnalysis.PiProofs.FiniteArchimedesBounds}
\lean{ComputableAnalysis.PiProofs.finiteArchimedesBounds}
\uses{def:archimedes-area-circumference-enclosures}
\leanok
For each positive polygon stage, the inner sector area is below the
corresponding chordal path estimate, and the tangent path estimate is below
the outer sector area after the usual normalization.  Equivalently, the area
and circumference intervals overlap at the same stage: the lower endpoint of
the area interval is below the upper endpoint of the circumference interval,
and the lower endpoint of the circumference interval is below the upper
endpoint of the area interval.
\end{theorem}

\begin{theorem}[Archimedes area-circumference agreement]
\label{thm:archimedes-area-circumference-agreement}
\lean{ComputableAnalysis.PiProofs.archimedesTheorem}
\uses{thm:finite-archimedes-bounds,def:raw-equiv}
\leanok
The polygonal area algorithm and the half-circumference algorithm compute
the same raw real.  In Lean this is the raw-equivalence theorem between
\(\code{ComputableAnalysis.piCircleArea}\) and
\(\code{ComputableAnalysis.piCircumference}\).
This is the Archimedes bridge between exhaustion by area and exhaustion by
length.
\end{theorem}

\begin{conjecture}[Geometric pi validity package]
\label{conj:geometric-pi-validity}
\lean{ComputableAnalysis.PiProofs.ValidityProofs}
\lean{ComputableAnalysis.PiProofs.DyadicCellRefinementBounds}
\lean{ComputableAnalysis.PiProofs.pointCross_split_difference_mul_den}
\lean{ComputableAnalysis.PiProofs.pointCross_split_le}
\lean{ComputableAnalysis.PiProofs.innerAreaCellSplitBound}
\lean{ComputableAnalysis.PiProofs.pointTangentCrossSum_formula_of_lt}
\lean{ComputableAnalysis.PiProofs.tangentWidth_split_difference_mul_den_of_nonneg}
\lean{ComputableAnalysis.PiProofs.tangentWidth_split_le}
\lean{ComputableAnalysis.PiProofs.outerAreaCellSplitBound}
\lean{ComputableAnalysis.PiProofs.innerEdgeCrosses_le_refinedInnerEdgeCrosses}
\lean{ComputableAnalysis.PiProofs.innerFanPerimeter_le_refinedInnerFanPerimeter}
\lean{ComputableAnalysis.PiProofs.innerQuarterArea_mono_of_inner_area_cell_bounds}
\lean{ComputableAnalysis.PiProofs.refinedOuterEdgeCrosses_le_outerEdgeCrosses}
\lean{ComputableAnalysis.PiProofs.refinedOuterFanPerimeter_le_outerFanPerimeter}
\lean{ComputableAnalysis.PiProofs.outerQuarterArea_antitone_of_outer_area_cell_bounds}
\lean{ComputableAnalysis.PiProofs.areaStageRefines_rationalCircle}
\lean{ComputableAnalysis.PiProofs.areaStepRefines_rationalCircle}
\lean{ComputableAnalysis.PiProofs.DyadicCellBoundsYieldStageRefinement}
\lean{ComputableAnalysis.PiProofs.rationalPointPathLength_cons_cons_lo}
\lean{ComputableAnalysis.PiProofs.rationalPointPathLength_cons_cons_hi}
\lean{ComputableAnalysis.PiProofs.innerPathLengthLo_le_refinedInnerPathLengthLo}
\lean{ComputableAnalysis.PiProofs.innerQuarterLength_mono_of_inner_length_cell_bounds}
\lean{ComputableAnalysis.PiProofs.refinedOuterPathLengthHi_le_outerPathLengthHi}
\lean{ComputableAnalysis.PiProofs.outerQuarterLength_antitone_of_outer_length_cell_bounds}
\lean{ComputableAnalysis.PiProofs.circumferenceStageRefines_of_length_cell_bounds}
\lean{ComputableAnalysis.PiProofs.dyadicCellBoundsYieldStageRefinement}
\lean{ComputableAnalysis.PiProofs.dyadicGeometricStageRefines_of_cellBounds}
\lean{ComputableAnalysis.PiProofs.geometricStepRemainders_of_dyadic_stage_refinement}
\lean{ComputableAnalysis.PiProofs.geometricStepRemainders_of_dyadic_cell_bounds}
\lean{ComputableAnalysis.PiProofs.areaValid_of_stepRefines_and_shrinking}
\lean{ComputableAnalysis.PiProofs.circumferenceValid_of_stepRefines_and_shrinking}
\uses{def:archimedes-area-circumference-enclosures,def:raw-real}
The area and circumference algorithms are valid raw reals: their intervals
are ordered, nested under dyadic refinement, and have widths tending to zero.
The nesting proof is organized locally.  A level \(n+1\) dyadic subdivision
splits each level \(n\) cell into two cells; the Lean interface isolates the
four local inequalities for inner area, outer area, inner length, and outer
length.  The two area inequalities are formalized.  The inner-area split is
the determinant identity for three ordered rational-circle points; the
outer-area split rewrites tangent widths as
\(2(v-u)/(1+uv)\) and clears denominators, leaving the nonnegative product
\((v-u)(w-v)(w-u)\).  Summing these local inequalities proves that the
inscribed quarter-area increases and the circumscribed quarter-area decreases
under dyadic refinement, giving the Lean theorem
\(\texttt{areaStepRefines\_rationalCircle}\).

For circumference, the same summation layer is now formalized with the
separate square-root precision schedule: coarse segment intervals are
computed at stage \((N+1)^3\), while the refined stage uses \((2N+1)^3\).
If the local inner and outer length split bounds hold at a stage, then the
inner quarter-length lower endpoint increases and the outer quarter-length
upper endpoint decreases at the refined stage.  Thus
\(\texttt{dyadicCellBoundsYieldStageRefinement}\) turns the four local cell
bounds into dyadic stage refinement.  What remains for full geometric
validity is to prove the two local length inequalities themselves, and then
give explicit width bounds for the area and length polygon gaps.
\end{conjecture}

\section{Trigonometric Functions}

The main function-level object is a raw function that sends an input angle to
a certified box around the corresponding point on the oriented unit circle:
\[
  \Theta \longmapsto C(\Theta)=(X(\Theta),Y(\Theta)).
\]
Cosine and sine are then coordinate projections:
\[
  \cos(\Theta)=X(\Theta),
  \qquad
  \sin(\Theta)=Y(\Theta).
\]

\begin{definition}[Rational-circle coordinates]
\label{def:rational-circle-sin-cos}
\lean{ComputableAnalysis.RationalCircle.Trigonometry.cos}
\lean{ComputableAnalysis.RationalCircle.Trigonometry.sin}
\uses{def:rational-circle-point}
For a rational slope coordinate \(u\), define the exact rational coordinates
\[
  c(u)=P(u)_x,\qquad s(u)=P(u)_y.
\]
Equivalently,
\[
  c(u)=\frac{1-u^2}{1+u^2},
  \qquad
  s(u)=\frac{2u}{1+u^2}.
\]
These are finite-stage coordinate functions, not yet the full angle-input
trigonometric functions.
\end{definition}

\begin{definition}[Other rational-circle coordinate functions]
\label{def:rational-circle-other-trig}
\lean{ComputableAnalysis.RationalCircle.Trigonometry.tan}
\lean{ComputableAnalysis.RationalCircle.Trigonometry.cot}
\lean{ComputableAnalysis.RationalCircle.Trigonometry.sec}
\lean{ComputableAnalysis.RationalCircle.Trigonometry.csc}
\uses{def:rational-circle-sin-cos}
On this rational circle layer, the quotient and reciprocal coordinate
functions are
\[
  t(u)=\frac{s(u)}{c(u)},\quad
  \operatorname{ct}(u)=\frac{c(u)}{s(u)},
\]
\[
  \operatorname{sc}(u)=\frac1{c(u)},\quad
  \operatorname{cs}(u)=\frac1{s(u)}.
\]
The full partial functions will later carry explicit denominator-apartness
domain data.
\end{definition}

\begin{definition}[Angle-to-circle-point raw function]
\label{def:geometric-point-functionraw}
\uses{def:dyadic-rational-circle-stages,def:archimedes-area-circumference-enclosures}
The central construction target is a raw function
\[
  C : \Q \to \hbox{boxes in }\Q^2
\]
whose value at a rational angle input \(\Theta\) encloses the point on the
oriented unit circle with signed sector area \(\Theta/2\), after period
reduction.

The intended algorithm is:
\begin{enumerate}
  \item reduce \(\Theta\) to a fundamental period interval;
  \item choose the quadrant and orientation;
  \item search over rational slope intervals using dyadic circle stages;
  \item use inner and outer rational sector areas to certify which arc
  contains the target;
  \item return the coordinate box of the certified arc, transformed back by
  the quadrant data.
\end{enumerate}
The algorithm may be refactored internally, but its public output should be a
single point raw function.
\end{definition}

\begin{definition}[Sine and cosine by coordinate projection]
\label{def:geometric-sin-cos-functionraw}
\lean{ComputableAnalysis.RationalCircle.GeometricTrig.cosFunctionRawOfPoint}
\lean{ComputableAnalysis.RationalCircle.GeometricTrig.sinFunctionRawOfPoint}
\lean{ComputableAnalysis.RationalCircle.GeometricTrig.cosFunctionRawOfPoint_compute_eq}
\lean{ComputableAnalysis.RationalCircle.GeometricTrig.sinFunctionRawOfPoint_compute_eq}
\uses{def:geometric-point-functionraw}
\leanok
Given a circle-point raw function \(C\), define
\[
  \cos_C(\Theta)=\operatorname{re}(C(\Theta)),
  \qquad
  \sin_C(\Theta)=\operatorname{im}(C(\Theta)).
\]
In Lean this is a projection from a complex-valued `FunctionRaw` on the real
axis to two `PartialRealFunRaw` coordinate algorithms.
\end{definition}

\begin{definition}[FunctionRaw construction package]
\label{def:geometric-trig-functionraw-construction}
\lean{ComputableAnalysis.RationalCircle.GeometricTrig.FunctionRawConstruction}
\lean{ComputableAnalysis.RationalCircle.GeometricTrig.FunctionRawConstruction.cosFunctionRaw}
\lean{ComputableAnalysis.RationalCircle.GeometricTrig.FunctionRawConstruction.sinFunctionRaw}
\uses{def:geometric-sin-cos-functionraw}
\leanok
The theorem-facing construction package records:
\[
  C,\quad \cos_C,\quad \sin_C,
\]
validity of the point and coordinate algorithms, the unit-circle law, and the
identity package to be proved for the projected functions.
\end{definition}

\section{Trigonometric Identities}

\begin{theorem}[Pythagorean identity on the rational circle]
\label{thm:rational-circle-pythagorean}
\lean{ComputableAnalysis.RationalCircle.Trigonometry.cos_sq_add_sin_sq}
\uses{def:rational-circle-sin-cos,thm:rational-circle-equation}
\leanok
For every rational slope coordinate \(u\),
\[
  c(u)^2+s(u)^2=1.
\]
\end{theorem}

\begin{theorem}[Basic rational-circle identity package]
\label{thm:rational-circle-basic-identities}
\lean{ComputableAnalysis.RationalCircle.Trigonometry.basicIdentityPackage}
\lean{ComputableAnalysis.RationalCircle.Trigonometry.composed_cos_eq}
\lean{ComputableAnalysis.RationalCircle.Trigonometry.composed_sin_eq}
\lean{ComputableAnalysis.RationalCircle.Trigonometry.double_cos_eq_sq_sub_sq}
\lean{ComputableAnalysis.RationalCircle.Trigonometry.double_sin_eq_two_mul}
\lean{ComputableAnalysis.RationalCircle.Trigonometry.one_add_tan_sq_eq_sec_sq}
\lean{ComputableAnalysis.RationalCircle.Trigonometry.one_add_cot_sq_eq_csc_sq}
\lean{ComputableAnalysis.RationalCircle.Trigonometry.cos_neg}
\lean{ComputableAnalysis.RationalCircle.Trigonometry.sin_neg}
\lean{ComputableAnalysis.RationalCircle.Trigonometry.tan_neg}
\uses{def:rational-circle-other-trig,thm:rational-circle-pythagorean}
\leanok
The rational-circle coordinates satisfy the quotient, reciprocal,
Pythagorean, parity, double-coordinate, and composition identities.  In
particular,
\[
  1+t(u)^2=\operatorname{sc}(u)^2\quad(c(u)\ne0),
  \qquad
  1+\operatorname{ct}(u)^2=\operatorname{cs}(u)^2\quad(s(u)\ne0),
\]
and
\[
  (P(u)P(v))_x=c(u)c(v)-s(u)s(v),
\]
\[
  (P(u)P(v))_y=c(u)s(v)+s(u)c(v).
\]
\end{theorem}

\begin{theorem}[Circle group law]
\label{thm:rational-circle-group-law}
\uses{def:rational-circle-point}
Let
\[
  P(u)\cdot P(v)
  =
  (P(u)_xP(v)_x-P(u)_yP(v)_y,\,
   P(u)_xP(v)_y+P(u)_yP(v)_x)
\]
be multiplication of unit-circle points.  In the stereographic chart, this
has a rational chart-change formula away from the missing point.  The missing
point is handled by the same quadrant or projective branch data used by the
geometric point algorithm.
\end{theorem}

\begin{theorem}[Projected addition formulas]
\label{thm:geometric-addition-formulas}
\uses{def:geometric-trig-functionraw-construction,thm:rational-circle-group-law}
Once the point raw function is proved compatible with circle multiplication,
the projected functions satisfy
\[
  \cos_C(a+b)=\cos_C(a)\cos_C(b)-\sin_C(a)\sin_C(b),
\]
\[
  \sin_C(a+b)=\sin_C(a)\cos_C(b)+\cos_C(a)\sin_C(b).
\]
The proof should first be finite and rational on circle boxes, then lifted to
the raw functions by validity and equivalence of the coordinate projections.
\end{theorem}

\begin{theorem}[Standard geometric trigonometric identities]
\label{thm:geometric-trig-identities}
\uses{thm:geometric-addition-formulas,thm:rational-circle-basic-identities}
The final identity package for this chapter contains:
\[
  \sin_C^2 x+\cos_C^2 x=1,\quad
  1+\tan_C^2 x=\sec_C^2 x,\quad
  1+\cot_C^2 x=\csc_C^2 x,
\]
the parity formulas, double-angle formulas, half-angle formulas, cofunction
formulas, and period formulas.  Each identity should be proved from the
circle point construction and its algebra, not by importing another
trigonometric representation.
\end{theorem}

\subsection{Formalization Order}

\begin{enumerate}
  \item Keep the exact rational-circle layer small and complete: points,
  dyadic refinement, circle multiplication, quotient functions, and the
  finite identity package.

  \item Complete the Archimedes method of exhaustion: polygonal area and
  circumference enclosures, same-stage overlap, validity, and the shared
  geometric pi real.

  \item Build the geometric point raw function as the public algorithm:
  sector-area bounds, dyadic bisection, period reduction, quadrant handling,
  validity, and shrinking widths.

  \item Define sine and cosine as coordinate projections of the point raw
  function.  Define the remaining trigonometric functions as partial
  quotient or reciprocal raw functions with explicit apartness domains.

  \item Prove the identities for the projected algorithms.  The main proof
  inputs are the unit-circle law, circle multiplication, and the correctness
  of period and quadrant reduction.

  \item Put comparison with any other representation in a later comparison
  chapter or file.  That later stage can import the other definitions it
  needs; this chapter should not depend on them.
\end{enumerate}
